\rtitle{Convergent Evolution of CAM Photosynthesis in Desert Plants: Evidence from Anatomy, Isotopes, and Gas Exchange}
\rsec{Introduction}
\clue{Crassulacean acid metabolism (CAM)} photosynthesis has \fact{evolved independently over 60 times across vascular plants, making it one of the most striking examples of evolutionary convergence on a complex trait} \rcite{1}. This pathway allows plants to open stomata at night when transpirational water loss is minimized, fixing CO\textsubscript{2} into malic acid for daytime use. The Agavoideae (\clue{agaves} and yuccas), Cactaceae (\clue{cacti}), and Euphorbiaceae (\clue{euphorbias}) represent three unrelated lineages that have independently evolved CAM and succulence as solutions to \clue{arid environments}. This report examines how \clue{anatomical cross-sections}, \clue{carbon isotope} \clue{ratios}, and \clue{nocturnal gas exchange} data reveal convergent evolutionary solutions across these lineages.

\rsec{CAM as a Convergent Adaptation to Aridity}
\rnote{Reference check. The epiphytes in the paragraph after Image~1 come from the figure caption on page [3] (an orchid, two bromeliads including Spanish moss); ``leafless photosynthetic roots'' is not in the page text or caption captured during the run.}
CAM photosynthesis is \fact{present in approximately 5\% of vascular plants}, distributed across \fact{ferns, gymnosperms, and angiosperms} \rcite{2}. The pathway involves \fact{nocturnal CO\textsubscript{2} fixation via phosphoenolpyruvate carboxylase (PEPC), storage of malic acid in vacuoles, and daytime decarboxylation} to supply CO\textsubscript{2} to Rubisco behind closed stomata \rcite{3}. This temporal separation of carbon acquisition and fixation dramatically improves water-use efficiency, enabling survival in deserts, epiphytic habitats, and other water-limited niches.

\rtrace{turn 8 \texttt{open(url)} \(\rightarrow\) turn 9 \texttt{open(id=page\_img\_215)}}{Source page: \emph{Prospects and perspectives: inferring physiological and regulatory targets for CAM from molecular and modelling approaches}; caption on the page: ``Diversity in succulent CAM lifeforms and a typical gas exchange profile. Terrestrial CAM plants: (A) massive stem succulent convergence in cacti (left) and euphorbia (right) species, with agave (centre)''\par\textbf{Why it was used.} The one photograph in the run with a cactus, a euphorbia and an agave side by side (panel A); panel I adds the four CAM phases over a day.}
\rimagebox{figures/showcase/cam/page_img_215__mcad142_fig1__8eaf997015.jpg}{Diversity of succulent CAM lifeforms including cacti, euphorbias, agaves, and epiphytic species, with a schematic of CAM phases showing nocturnal stomatal opening and daytime CO\textsubscript{2} release}{\rbox{0.0}{0.71}{0.51}{1.0}{}}
The morphological diversity of CAM plants is remarkable. Terrestrial representatives include \fact{massive stem succulents like cacti and euphorbias, leaf succulents like agaves, and more herbaceous forms like Crassulaceae} \rcite{3}.\rimgref{1} Epiphytic CAM plants include orchids, bromeliads, and Spanish moss. This diversity demonstrates that CAM can be integrated into many body plans, from arborescent forms to leafless photosynthetic roots.

\rsec{Anatomical Convergence: Leaf Structure and Succulence}
\rnote{Reference check. The figure on page [4] labels the vascular bundles, mesophyll and intercellular airspace of the three \emph{Yucca} leaves; it does not say that the CAM leaves are thicker or that their airspaces are large.}
Anatomical cross-sections reveal convergent structural adaptations for CAM function. The Agavoideae subfamily provides an excellent model for studying CAM evolution because it \fact{contains species spanning the C3--CAM continuum} \rcite{4}.

\rtrace{turn 2 \texttt{open(url)} \(\rightarrow\) turn 3 \texttt{open(id=page\_img\_23)}}{Source page: \emph{The Agavoideae: an emergent model clade for CAM evolutionary biology}; caption on the page: ``(A) Anatomical cross-sections of leaves from Yucca aloifolia (top, CAM), Y. gloriosa (middle, C3 + CAM) and Y. filamentosa (bottom, C3). VB: vascular bundle, M: undifferentiated mesophyll, IAS: intercellular airspace''\par\textbf{Why it was used.} The same figure as in the Concise report: a CAM, an intermediate and a C3 leaf from one genus. The text relies on the three sections in the left column.}
\rimagebox{figures/showcase/cam/page_img_23__mcad062_fig3__9018d4c929.jpg}{Anatomical cross-sections of Yucca leaves showing CAM (Y. aloifolia), C3+CAM intermediate (Y. gloriosa), and C3 (Y. filamentosa) phenotypes with labeled vascular bundles, mesophyll, and intercellular air spaces}{\rbox{0.0}{0.0}{0.31}{1.0}{}}
In \emph{Yucca}, CAM species like \emph{Y. aloifolia} exhibit thick leaves with large intercellular air spaces and \fact{undifferentiated mesophyll}, while C3 species like \emph{Y. filamentosa} show more typical leaf anatomy \rcite{4}.\rimgref{2} The hybrid \emph{Y. gloriosa} \fact{displays intermediate anatomy and physiology, making it a valuable system for studying CAM evolution} \rcite{8}.\rimgref{3}

\rtrace{turn 8 \texttt{open(url)} \(\rightarrow\) turn 9 \texttt{open(id=page\_img\_240)}}{Source page: \emph{Gas exchange and leaf anatomy of a C3--CAM hybrid, Yucca gloriosa (Asparagaceae)}; caption on the page: ``Cross sections of Y. aloifolia (A and B), Y. gloriosa (C and D), and Y. filamentosa (E and F) at 5\(\times\) (A, C, E) and 10\(\times\) (B, D, F) magnification.''\par\textbf{Why it was used.} Shows the hybrid's leaf between its two parents at two magnifications, for the sentence on \emph{Y. gloriosa}.}
\rimagebox{figures/showcase/cam/page_img_240__exbotj_erv536_f0003__314ac5b4c8.jpg}{Cross-sections of Y. aloifolia (CAM), Y. gloriosa (C3+CAM), and Y. filamentosa (C3) at 5\(\times\) and 10\(\times\) magnification showing progressive anatomical differences}{}
The \emph{Euphorbia milii} leaf cross-section reveals another convergent anatomical solution. The leaf contains \fact{palisade parenchyma, spongy parenchyma, and a bundle sheath with chloroplasts, suggesting possible C\textsubscript{2} photosynthesis} alongside \fact{weak CAM} \rcite{5}.\rimgref{4}

\rtrace{turn 5 \texttt{open(url)} \(\rightarrow\) turn 6 \texttt{open(id=page\_img\_139)}}{Source page: \emph{Crassulacean acid metabolism-cycling in Euphorbia milii}; caption on the page: ``Cross-sections of the leaf of E. milii. UE, upper epidermis; PP, palisade parenchyma; VB, vascular bundle; BS, bundle sheath; SP, spongy parenchyma; LE, lower epidermis. Arrowheads point at chloroplasts.''\par\textbf{Why it was used.} The run's only cross-section from a euphorbia, one of the three lineages in the request.}
\rimagebox{figures/showcase/cam/page_img_139__plt01401__9986f5efbb.jpg}{Cross-section of E. milii leaf showing upper epidermis, palisade parenchyma, vascular bundle, bundle sheath, spongy parenchyma, and lower epidermis}{}
\fact{Succulence}---fleshy tissue with high water content---\fact{is hypothesized to be essential for CAM function because it allows storage of large amounts of assimilated carbon as malate and assists in trapping CO\textsubscript{2} released during the day} \rcite{1}.

\rsec{Carbon Isotope Ratios: Quantifying CAM Contribution}
\rnote{Reference check. Page [6] states that PEP carboxylase discriminates less against \textsuperscript{13}C and that strong CAM plants resemble C4 plants, but the ranges in parentheses are not in the text captured from it. The peaks at \(-\)13\textperthousand{} and \(-\)26\textperthousand{} that the report gives for the 506-species figure come from a bromeliad survey described elsewhere on the page; the figure itself peaks near \(-\)14\textperthousand{} and \(-\)28\textperthousand{}. Page [2] does not say which Aizoaceae subfamilies are CAM.}
Carbon isotope ratios (\(\delta\)\textsuperscript{1}\textsuperscript{3}C) provide a powerful tool for identifying CAM plants and quantifying the relative contribution of nocturnal versus diurnal carbon fixation. The enzyme \fact{PEP carboxylase discriminates less against \textsuperscript{1}\textsuperscript{3}C than Rubisco, so plants with strong CAM exhibit \(\delta\)\textsuperscript{1}\textsuperscript{3}C values similar to C4 plants} (around --10\textperthousand{} to --15\textperthousand{}), while C3 plants show more negative values (around --25\textperthousand{} to --30\textperthousand{}) \rcite{6}.

\rtrace{turn 2 \texttt{open(url)} \(\rightarrow\) turn 3 \texttt{open(id=page\_img\_78)}}{Source page: \emph{How Closely Do the \(\delta^{13}\)C Values of Crassulacean Acid Metabolism Plants Reflect the Proportion of CO\textsubscript{2} Fixed during Day and Night?}; caption on the page: ``Relationships between the cumulative net carbon gain in the light and the dark and the \(\delta^{13}\)C values \ldots{} Extrapolation of the fitted regression \ldots{} in B indicates \(\delta^{13}\)C values of \(-\)8.7\textperthousand{} for 100\% dark CO\textsubscript{2} uptake and \(-\)26.9\textperthousand{} for 100\% CO\textsubscript{2} uptake during the light.''\par\textbf{Why it was used.} The source of the two end-point values in the next paragraph, which come from panel B.}
\rimagebox{figures/showcase/cam/page_img_78__pp0827338006__f54f4be844.jpg}{Relationships between cumulative net carbon gain in light and dark and \(\delta\)\textsuperscript{1}\textsuperscript{3}C values, showing linear correlations for CAM species}{\rbox{0.06}{0.40}{0.95}{0.70}{}}
Experimental studies demonstrate that \fact{\(\delta\)\textsuperscript{1}\textsuperscript{3}C values correlate linearly with the proportion of carbon gained in the dark versus light}. For CAM species, extrapolation indicates \(\delta\)\textsuperscript{1}\textsuperscript{3}C values of approximately \fact{--8.7\textperthousand{} for 100\% dark CO\textsubscript{2} uptake and --26.9\textperthousand{} for 100\% light CO\textsubscript{2} uptake} \rcite{6}.\rimgref{5} This relationship allows researchers to estimate the CAM contribution in natural populations.

\rtrace{turn 2 \texttt{open(url)} \(\rightarrow\) turn 3 \texttt{open(id=page\_img\_79)}}{Source page: \emph{How Closely Do the \(\delta^{13}\)C Values of Crassulacean Acid Metabolism Plants Reflect the Proportion of CO\textsubscript{2} Fixed during Day and Night?}; caption on the page: ``\(\delta^{13}\)C values of 506 species from the Asclepiadaceae, Bromeliaceae, Crassulaceae, Cucurbitaceae, Didiereaceae, Euphorbiaceae, Orchidaceae, Polypodiaceae, and Vittariaceae plant families known to contain both C3 and CAM members.''\par\textbf{Why it was used.} A two-peaked histogram across nine families, including the Euphorbiaceae, for the sentence on the C3--CAM split.}
\rimagebox{figures/showcase/cam/page_img_79__pp0827338007__e697edd027.jpg}{\(\delta\)\textsuperscript{1}\textsuperscript{3}C values of 506 species from multiple plant families showing bimodal distribution with peaks around --13\textperthousand{} (strong CAM) and --26\textperthousand{} (C3)}{}
The bimodal distribution of \(\delta\)\textsuperscript{1}\textsuperscript{3}C values across plant families reveals the evolutionary transition between C3 and CAM photosynthesis.\rimgref{6} In the Aizoaceae, for example, different subfamilies show distinct \(\delta\)\textsuperscript{1}\textsuperscript{3}C distributions reflecting their CAM status \rcite{2}.\rimgref{7}

\rtrace{turn 5 \texttt{open(url)} \(\rightarrow\) turn 6 \texttt{open(id=page\_img\_116)}}{Source page: \emph{The diverse diaspora of CAM: a pole-to-pole sketch}; caption on the page: ``Values of \(\delta^{13}\)C (per mille) of taxa from the five subfamilies of the Aizoaceae. Note that frequency values include values from different species and multiple measurements of the same species.''\par\textbf{Why it was used.} One family broken down by subfamily, for the Aizoaceae sentence.}
\rimagebox{figures/showcase/cam/page_img_116__mcad067_fig5__6e4363d21a.jpg}{\(\delta\)\textsuperscript{1}\textsuperscript{3}C values of taxa from five Aizoaceae subfamilies showing variation from C3 to strong CAM}{}
In \emph{Euphorbia milii}, \fact{\(\delta\)\textsuperscript{1}\textsuperscript{3}C values of --25.2\textperthousand{} in leaves and --24.7\textperthousand{} in stems} indicate \fact{weak CAM or CAM-cycling, intermediate between C3 and constitutive CAM plants} \rcite{5}. This intermediate value \fact{results from recycling of nocturnal respiratory CO\textsubscript{2}} rather than direct atmospheric fixation.

\rsec{Nocturnal Gas Exchange: The CAM Rhythm}
Continuous gas exchange measurements reveal the characteristic CAM rhythm of nocturnal CO\textsubscript{2} uptake and daytime release. In \emph{Agave angustifolia}, \fact{CO\textsubscript{2} uptake in the dark contributes about 75\% of daily carbon gain, and this pattern is highly conserved across environmental conditions} \rcite{7}.\rimgref{8}

\rtrace{turn 2 \texttt{open(url)} \(\rightarrow\) turn 3 \texttt{open(id=page\_img\_45)}}{Source page: \emph{Nocturnal versus diurnal CO\textsubscript{2} uptake: how flexible is Agave angustifolia?}; caption on the page: ``Net CO\textsubscript{2} balance of the shoot of a young individual of A. angustifolia during 281 consecutive 12h light/12h dark cycles. Carbon balances are shown for the light and dark periods.''\par\textbf{Why it was used.} The run's only long gas-exchange record from an agave: nine months of day and night carbon balance under changing light and watering.}
\rimagebox{figures/showcase/cam/page_img_45__exbotj_eru097_f0001__bfdd7f01a3.jpg}{Net CO\textsubscript{2} balance of young A. angustifolia during 281 consecutive 12h light/12h dark cycles showing consistent nocturnal CO\textsubscript{2} uptake}{}
The nocturnal CO\textsubscript{2} uptake in CAM plants is not simply a matter of stomatal opening. Research on \emph{Clusia} species demonstrates that \fact{CAM requires elevated nocturnal respiratory rates to power vacuolar malic acid accumulation} \rcite{9}.

\rtrace{turn 8 \texttt{open(url)} \(\rightarrow\) turn 9 \texttt{open(id=page\_img\_264)}}{Source page: \emph{Elevated nocturnal respiratory rates in the mitochondria of CAM plants: current knowledge and unanswered questions}; caption on the page: ``Schematic of major metabolite fluxes involved during the nocturnal phase of CAM and their interaction with the respiratory TCA cycle.''\par\textbf{Why it was used.} Shows where respiration feeds malate synthesis at night, for the paragraph on the metabolic cost of CAM.}
\rimagebox{figures/showcase/cam/page_img_264__mcad119_fig1__ce0ef55474.jpg}{Schematic of metabolite fluxes during nocturnal CAM phase showing interaction between mitochondrial respiration, TCA cycle, and malate synthesis}{}
The metabolic demands of CAM are substantial. Nocturnal respiration provides the ATP and reducing power needed for PEP regeneration and malate synthesis.\rimgref{9} \fact{Flux balance analysis models predict that CAM requires elevated nocturnal respiratory rates compared to C3 species} \rcite{9}.

\rtrace{turn 8 \texttt{open(url)} \(\rightarrow\) turn 9 \texttt{open(id=page\_img\_266)}}{Source page: \emph{Elevated nocturnal respiratory rates in the mitochondria of CAM plants: current knowledge and unanswered questions}; caption on the page: ``Clusia pratensis (grey dots) and C. tocuchensis (white dots) exhibit facultative CAM and obligate C3 phenotypes, respectively.''\par\textbf{Why it was used.} Pairs a facultative CAM species with a C3 relative under watering and drought, for the sentence on drought-induced CAM.}
\rimagebox{figures/showcase/cam/page_img_266__mcad119_fig3__d8fa6cd60f.jpg}{Diel photosynthetic assimilation rates in well-watered and drought-treated Clusia pratensis (facultative CAM) and C. tocuchensis (obligate C3)}{}
In \fact{facultative CAM} species like \emph{Clusia pratensis}, drought induces CAM, with nocturnal CO\textsubscript{2} uptake becoming prominent under water stress \rcite{9}.\rimgref{10} This plasticity allows plants to optimize carbon gain across varying water availability.

\rsec{Facultative CAM and Environmental Plasticity}
\rnote{Figure check. In the \emph{Lecanopteris} record, drying lowers daytime uptake and adds a small night-time uptake; night uptake never becomes the larger part. In the \emph{E. milii} time course, total acid content rises under drought, but nocturnal accumulation (panel C) does not: at pH 7.0 it stays level, and at pH 8.4 it peaks at day 5 and falls to zero by day 16.}
Not all CAM plants are obligate. Some species, like the ant-fern \emph{Lecanopteris sinuosa}, exhibit \fact{facultative CAM} that can be induced by drought \rcite{2}.

\rtrace{turn 5 \texttt{open(url)} \(\rightarrow\) turn 6 \texttt{open(id=page\_img\_112)}}{Source page: \emph{The diverse diaspora of CAM: a pole-to-pole sketch}; caption on the page: ``Demonstration of facultative CAM in the ant-fern Lecanopteris sinuosa (Polypodiaceae) during 20 days of net CO\textsubscript{2} exchange \ldots{} Net CO\textsubscript{2} exchange during the light is shown as yellow, whereas net CO\textsubscript{2} exchange during the dark is shown as blue.''\par\textbf{Why it was used.} A watering--drying--rewatering cycle in a fern, showing that CAM can switch on and off within weeks.}
\rimagebox{figures/showcase/cam/page_img_112__mcad067_fig1__34fcbf4429.jpg}{Demonstration of facultative CAM in Lecanopteris sinuosa during watering-drying-rewatering cycles, with net CO\textsubscript{2} exchange showing nocturnal uptake during drought}{}
In \emph{Lecanopteris sinuosa}, net CO\textsubscript{2} exchange shifts from predominantly diurnal (yellow) to nocturnal (blue) uptake during drought periods, demonstrating the plasticity of CAM expression \rcite{2}.\rimgref{11}

Similarly, \emph{Euphorbia milii} shows \fact{CAM-cycling} rather than strong CAM, with \fact{nocturnal H\textsuperscript{+} accumulation attributable to both malate and citrate} \rcite{5}.

\rtrace{turn 5 \texttt{open(url)} \(\rightarrow\) turn 6 \texttt{open(id=page\_img\_140)}}{Source page: \emph{Crassulacean acid metabolism-cycling in Euphorbia milii}; caption on the page: ``Time course of changes with drought in leaves of E. milii in (A) H\textsuperscript{+} content titrated to pH 7.0 \ldots{} (C) nocturnal H\textsuperscript{+} accumulation \ldots{} and (D) dawn leaf FM per area (circles) and chlorophyll content (triangles).''\par\textbf{Why it was used.} The drought time course for the euphorbia, for the sentence on acid accumulation under drought.}
\rimagebox{figures/showcase/cam/page_img_140__plt01402__c3d4baca77.jpg}{Time course of changes with drought in E. milii leaves showing H\textsuperscript{+} content at pH 7.0 and 8.4, nocturnal H\textsuperscript{+} accumulation, and chlorophyll content}{\rbox{0.08}{0.27}{0.88}{0.49}{}}
Under drought, \emph{E. milii} shows increased nocturnal H\textsuperscript{+} accumulation, but the pattern differs from strong CAM, with citrate contributing significantly to acidification \rcite{5}.\rimgref{12}

\rsec{Phylogenetic Distribution and Evolutionary Origins}
The phylogenetic distribution of CAM reveals its polyphyletic nature. In the Portulacineae plus Molluginaceae, CAM has evolved multiple times independently \rcite{2}.

\rtrace{turn 5 \texttt{open(url)} \(\rightarrow\) turn 6 \texttt{open(id=page\_img\_117)}}{Source page: \emph{The diverse diaspora of CAM: a pole-to-pole sketch}; caption on the page: ``Presence of C3, C3 + CAM, strong CAM, C4 + CAM, C4 and C4 intermediate photosynthesis (indicated by green ticks) superimposed upon phylogenies of the Portulacineae plus Molluginaceae (upper panel) and the Montiaceae (lower panel).''\par\textbf{Why it was used.} Places the photosynthetic types on a phylogeny that includes the Cactaceae, for the section on how often CAM arose.}
\rimagebox{figures/showcase/cam/page_img_117__mcad067_fig6__db4923e3f4.jpg}{Phylogeny of Portulacineae plus Molluginaceae showing distribution of C3, C3+CAM, CAM, C4+CAM, C4, and C4-intermediate photosynthesis}{}
Within the Montiaceae, \fact{different genera show different photosynthetic variants, with some species exhibiting C3+CAM while others are strictly C3} \rcite{2}.\rimgref{13} This pattern suggests that CAM can evolve within lineages that are otherwise C3, and that intermediate forms like C3+CAM are common stepping stones.

\rsec{Conclusions}
The evidence from anatomical cross-sections, carbon isotope ratios, and nocturnal gas exchange demonstrates that CAM photosynthesis represents a remarkable case of convergent evolution. Unrelated desert plant lineages---agaves, cacti, and euphorbias---have independently evolved similar solutions to aridity: nocturnal CO\textsubscript{2} fixation, succulent tissue for malate storage, and temporal separation of carbon acquisition from fixation.

Anatomical studies reveal convergent leaf structures with large vacuoles for acid storage and modified mesophyll organization. Carbon isotope ratios provide quantitative measures of CAM contribution, revealing a continuum from C3 through C3+CAM to strong CAM. Nocturnal gas exchange measurements confirm the characteristic CAM rhythm and demonstrate the metabolic costs of this pathway, including elevated nocturnal respiration.

The diversity of CAM expression---from obligate CAM in agaves to facultative CAM in ferns and CAM-cycling in \emph{Euphorbia}---demonstrates that this pathway can be integrated into many physiological and ecological contexts. As climate change increases water stress globally, understanding the convergent evolution of CAM may inform efforts to engineer drought tolerance into crops.

\rsec{References}
\rcite{1} Crassulacean acid metabolism (CAM) at the crossroads: a special issue to honour 50 years of CAM research by Klaus Winter\rurl{https://pmc.ncbi.nlm.nih.gov/articles/PMC10799977/}

\rcite{2} The diverse diaspora of CAM: a pole-to-pole sketch\rurl{https://pmc.ncbi.nlm.nih.gov/articles/PMC10800000/}

\rcite{3} Prospects and perspectives: inferring physiological and regulatory targets for CAM from molecular and modelling approaches\rurl{https://pmc.ncbi.nlm.nih.gov/articles/PMC10799989/}

\rcite{4} The Agavoideae: an emergent model clade for CAM evolutionary biology\rurl{https://pmc.ncbi.nlm.nih.gov/articles/PMC10799990/}

\rcite{5} Crassulacean acid metabolism-cycling in Euphorbia milii\rurl{https://pmc.ncbi.nlm.nih.gov/articles/PMC3628315/}

\rcite{6} How Closely Do the \(\delta\)\textsuperscript{1}\textsuperscript{3}C Values of Crassulacean Acid Metabolism Plants Reflect the Proportion of CO\textsubscript{2} Fixed during Day and Night?\rurl{https://pmc.ncbi.nlm.nih.gov/articles/PMC166772/}

\rcite{7} Nocturnal versus diurnal CO\textsubscript{2} uptake: how flexible is Agave angustifolia?\rurl{https://pmc.ncbi.nlm.nih.gov/articles/PMC4085958/}

\rcite{8} Gas exchange and leaf anatomy of a C3--CAM hybrid, Yucca gloriosa (Asparagaceae)\rurl{https://pmc.ncbi.nlm.nih.gov/articles/PMC4762382/}

\rcite{9} Elevated nocturnal respiratory rates in the mitochondria of CAM plants: current knowledge and unanswered questions\rurl{https://pmc.ncbi.nlm.nih.gov/articles/PMC10799998/}

\rimgsources{\rimgsrc{Image~1}{Prospects and perspectives: inferring physiological and regulatory targets for CAM from molecular and modelling approaches}{https://pmc.ncbi.nlm.nih.gov/articles/PMC10799989/}
\rimgsrc{Image~2}{The Agavoideae: an emergent model clade for CAM evolutionary biology}{https://pmc.ncbi.nlm.nih.gov/articles/PMC10799990/}
\rimgsrc{Image~3}{Gas exchange and leaf anatomy of a C3--CAM hybrid, Yucca gloriosa (Asparagaceae)}{https://pmc.ncbi.nlm.nih.gov/articles/PMC4762382/}
\rimgsrc{Images~4, 12}{Crassulacean acid metabolism-cycling in Euphorbia milii}{https://pmc.ncbi.nlm.nih.gov/articles/PMC3628315/}
\rimgsrc{Images~5, 6}{How Closely Do the \(\delta^{13}\)C Values of Crassulacean Acid Metabolism Plants Reflect the Proportion of CO\textsubscript{2} Fixed during Day and Night?}{https://pmc.ncbi.nlm.nih.gov/articles/PMC166772/}
\rimgsrc{Images~7, 11, 13}{The diverse diaspora of CAM: a pole-to-pole sketch}{https://pmc.ncbi.nlm.nih.gov/articles/PMC10800000/}
\rimgsrc{Image~8}{Nocturnal versus diurnal CO\textsubscript{2} uptake: how flexible is Agave angustifolia?}{https://pmc.ncbi.nlm.nih.gov/articles/PMC4085958/}
\rimgsrc{Images~9, 10}{Elevated nocturnal respiratory rates in the mitochondria of CAM plants: current knowledge and unanswered questions}{https://pmc.ncbi.nlm.nih.gov/articles/PMC10799998/}}
\rgate{\item Submission 1 (turn 10): {\color{tok}\textbf{OK}} --- 7,377 chars, 13 images, 9 sources}
